% ===== PRÉAMBULE CANONIQUE — à coller VERBATIM en tête de CHAQUE .tex =====
\documentclass[11pt,a4paper]{article}
\usepackage[utf8]{inputenc}\usepackage[T1]{fontenc}\usepackage{lmodern}
\usepackage[french]{babel}
\usepackage{amsmath,amssymb,mathtools}\usepackage{icomma}
\usepackage{siunitx}\sisetup{locale=FR}  % \qty{}{}, \si{}, \ang{}
\usepackage{xcolor}\usepackage{colortbl}
\usepackage{tikz}\usepackage{pgfplots}\pgfplotsset{compat=1.18}
\usepackage[siunitx,europeanresistors]{circuitikz} % circuits (élec.)
\usepackage[most]{tcolorbox}\usepackage{enumitem}\usepackage{array,booktabs}
\usepackage[a4paper,margin=2.2cm]{geometry}
\usetikzlibrary{arrows.meta,calc,angles,quotes,positioning,patterns,%
  backgrounds,shapes.geometric,decorations.pathmorphing,decorations.markings,intersections}
\definecolor{primary}{RGB}{222,49,99}\definecolor{primarydark}{RGB}{150,22,55}
\definecolor{defcol}{RGB}{0,114,178}\definecolor{methcol}{RGB}{52,92,148}
\definecolor{excol}{RGB}{0,135,90}\definecolor{attncol}{RGB}{200,40,55}
\tcbset{boxbase/.style={enhanced,breakable,boxrule=0.9pt,arc=2.5pt,
  left=3mm,right=3mm,top=2.3mm,bottom=2.3mm,fonttitle=\bfseries,coltitle=white}}
% --- boîtes TUTORIEL (noms EXACTS, ne pas renommer) ---
\newtcolorbox{cle}[1][]{boxbase,colback=defcol!6,colframe=defcol,
  title={\textsc{À retenir}\ifx\relax#1\relax\relax\else\textmd{ : \itshape #1}\fi}}
\newtcolorbox{methode}[1][]{boxbase,colback=methcol!7,colframe=methcol,sharp corners,
  title={\textsc{Méthode}\ifx\relax#1\relax\relax\else\textmd{ --- \itshape #1}\fi}}
\newtcolorbox{exemplebox}[1][]{boxbase,colback=excol!5,colframe=excol,
  title={\textsc{Exemple}\ifx\relax#1\relax\relax\else\textmd{ : \itshape #1}\fi}}
\newtcolorbox{piege}[1][]{boxbase,colback=attncol!6,colframe=attncol,
  title={\textsc{Piège}\ifx\relax#1\relax\relax\else\textmd{ : \itshape #1}\fi}}
\newtcolorbox{remarque}[1][]{boxbase,colback=black!4,colframe=black!45,
  title={\textsc{Remarque}\ifx\relax#1\relax\relax\else\textmd{ : \itshape #1}\fi}}
% --- boîtes EXERCICE / CORRIGÉ (noms EXACTS, auto-comptées) ---
\newtcolorbox[auto counter]{exo}[1][]{boxbase,colback=primary!4,colframe=primary,
  title={\textsc{Exercice}~\thetcbcounter\ifx\relax#1\relax\relax\else\textmd{ --- \itshape #1}\fi}}
\newtcolorbox[auto counter]{corrige}[1][]{boxbase,colback=excol!4,colframe=excol,
  title={\textsc{Corrigé}~\thetcbcounter\ifx\relax#1\relax\relax\else\textmd{ --- \itshape #1}\fi}}
\newcommand{\vv}[1]{\vec{#1}}\newcommand{\ds}{\displaystyle}
\newcommand{\R}{\mathbb{R}}\newcommand{\N}{\mathbb{N}}\newcommand{\Z}{\mathbb{Z}}
% =========================================================================
\begin{document}
\sisetup{per-mode=symbol}

\begin{corrige}[classer trois descentes]
Sans frottement, $a=g\sin\alpha$ : l'accélération \emph{ne dépend que de l'angle}, pas de la masse.
Comme $\sin\ang{20}<\sin\ang{30}<\sin\ang{50}$, on a
\[ a_A<a_B<a_C. \]
Les masses ($6$, $2$, $4$ kg) sont des \emph{données-pièges} : elles n'interviennent pas.
\end{corrige}

\begin{corrige}[le cycliste : équilibre puis accélération]
\begin{enumerate}[label=\alph*)]
  \item Vitesse constante $\Rightarrow a=0$ : le frottement équilibre la composante motrice,
        $F_f=G\sin\ang{4}=802\times\num{0,0698}\approx\qty{55,9}{\newton}$.
  \item Masse : $m=\dfrac{G}{g}=\dfrac{802}{10}=\qty{80,2}{\kilogram}$. Composante motrice
        $G_x=G\sin\ang{30}=802\times\num{0,5}=\qty{401}{\newton}$. Selon $Ox$ :
        \[ a=\frac{G_x-F_f}{m}=\frac{401-20}{\num{80,2}}=\frac{381}{\num{80,2}}\approx\qty{4,75}{\metre\per\second\squared}. \]
\end{enumerate}
\end{corrige}

\begin{corrige}[frottement statique : la force qui s'ajuste]
Sur l'horizontale, $N=mg=\qty{100}{\newton}$, donc $F_{f,\max}=\mu_s N=\num{0,5}\times100=\qty{50}{\newton}$.
\begin{enumerate}[label=\alph*)]
  \item Immobile $\Rightarrow a=0$ : le frottement statique \emph{égale} la traction, $F_f=\qty{40}{\newton}$.
  \item $40<F_{f,\max}=50$ : la caisse peut rester immobile. \checkmark
  \item À \qty{45}{\newton} et toujours immobile, $F_f=\qty{45}{\newton}$ (le frottement s'ajuste). Dès que la
        traction dépasse $F_{f,\max}=\qty{50}{\newton}$, l'équilibre est rompu et la caisse se met à glisser.
\end{enumerate}
\end{corrige}

\begin{corrige}[coefficient de frottement par un MRUA]
Sur l'horizontale, $N=mg$ et le frottement est $F_f=\mu_c mg$. Selon la direction du mouvement :
\[ F-\mu_c mg=ma \;\Rightarrow\; \mu_c=\frac{F-ma}{mg}=\frac{400-50\times2}{50\times10}
   =\frac{300}{500}=\num{0,6}. \]
\end{corrige}

\begin{corrige}[vitesse au bas d'une pente]
\begin{enumerate}[label=\alph*)]
  \item $a=g\sin\alpha=10\times\num{0,5}=\qty{5}{\metre\per\second\squared}$.
  \item $v=\sqrt{2a\,d}=\sqrt{2\times5\times10}=\sqrt{100}=\qty{10}{\metre\per\second}$.
\end{enumerate}
\end{corrige}
\end{document}
